% Part: incompleteness
% Chapter: theories-computability
% Section: first-incompleteness

\documentclass[../../../include/open-logic-section]{subfiles}

\begin{document}

\olfileid[hi]{inc}{tcp}{inc}

\olsection{$\Th{Q}$ का कोई पूर्ण, संगत, !!^{axiomatizable} विस्तार नहीं है}

\begin{thm}
\ollabel{thm:first-incompleteness}
$\Th{Q}$ का कोई पूर्ण, संगत, !!{axiomatizable} विस्तार नहीं है।
\end{thm}

\begin{proof}
हम पहले ही जानते हैं कि $\Th{Q}$ का कोई संगत, निर्णेय विस्तार नहीं है।
किंतु यदि $\Th{T}$ पूर्ण और !!{axiomatized} है, तो वह निर्णेय है।
\end{proof}

\begin{explain}
यह प्रमेय ग्योडेल (G\"odel) के 1931 के प्रथम अपूर्णता-प्रमेय के मूल
प्रतिपादन से बहुत दूर नहीं है। अधिक आधुनिक शब्दावली के अतिरिक्त मुख्य
अंतर ये हैं: ग्योडेल (G\"odel) के यहाँ ``संगत'' के स्थान पर ``$\omega$-संगत'' है;
और वह पूरी व्यापकता में ``!!{axiomatizable}'' नहीं कह सकते थे, क्योंकि उस
समय संगणनीयता की औपचारिक धारणा उपलब्ध नहीं थी। (अगले दशक में संगणनीयता के
औपचारिक प्रतिरूप विकसित किए गए---इनमें ग्योडेल (G\"odel) भी सहभागी थे---और इसका एक
बड़ा उद्देश्य उन प्रकार के सिद्धांतों का अभिलक्षण करना था जो ग्योडेलीय (G\"odelian)
परिघटना के अधीन होते हैं।)

प्रमेय कहता है कि आपको सब कुछ एक साथ नहीं मिल सकता: पूर्णता, संगतता और
!!{axiomatizability}। किंतु इनमें से किसी एक को छोड़ दें तो अन्य दो मिल
सकते हैं: $\Th{Q}$ संगत और संगणनीयतः अभिगृहीतीकृत है, किंतु पूर्ण नहीं;
असंगत सिद्धांत पूर्ण और संगणनीयतः अभिगृहीतीकृत (उदाहरणतः
$\{ \eq/[0][0] \}$ द्वारा) है, किंतु संगत नहीं; और अंकगणित के सत्य वाक्यों
का समुच्चय पूर्ण तथा संगत है, किंतु संगणनीयतः अभिगृहीतीकृत नहीं।
\end{explain}
\end{document}
